% Section 4: Implementation

Empirica is an open-source framework implementing the epistemic vector theory. It provides infrastructure for AI self-assessment, epistemic gating, and calibration feedback loops.

\subsection{Architecture Overview}

The framework consists of eight core components:

\begin{enumerate}
    \item \textbf{Sessions Database}: Persistent storage for session state, including AI identity, timestamps, and metadata.
    \item \textbf{Cascades}: Complete PREFLIGHT $\rightarrow$ CHECK $\rightarrow$ POSTFLIGHT $\rightarrow$ POST-TEST loops that capture epistemic state evolution and ground self-assessment in objective evidence.
    \item \textbf{Bayesian Belief Manager}: Maintains dual-track calibration state---self-referential corrections from learning deltas and grounded corrections from post-test evidence.
    \item \textbf{Sentinel Protocol}: The epistemic gate that enforces noetic-before-praxic via capability-based access control.
    \item \textbf{Post-Test Verification}: Collects objective evidence after task completion to ground self-assessment in measurable outcomes.
    \item \textbf{Epistemic Intent Layer}: Tracks assumptions (unverified beliefs with urgency decay), decisions (audit-trail choice points), and intent edges (noetic-to-praxic provenance).
    \item \textbf{Procedural Knowledge (Lessons)}: Executable procedures extracted from successful patterns, subject to the cognitive immune system.
    \item \textbf{Semantic Memory}: Four-layer architecture (HOT/WARM/SEARCH/COLD) with Qdrant-backed retrieval for contextual injection at PREFLIGHT.
\end{enumerate}

\subsection{The CASCADE Workflow}

CASCADE (Calibrated Assessment of State, Context, and Domain Epistemics) is the core epistemic loop:

\begin{figure}[h]
\centering
\begin{tikzpicture}[
    node distance=1cm,
    phase/.style={rectangle, draw, rounded corners, minimum width=2.8cm, minimum height=0.8cm, fill=blue!10, font=\small\bfseries},
    decision/.style={diamond, draw, aspect=2.5, fill=yellow!20, font=\small},
    action/.style={rectangle, draw, rounded corners, minimum width=2.2cm, minimum height=0.7cm, font=\small},
    update/.style={rectangle, draw, rounded corners, minimum width=2.5cm, minimum height=0.8cm, fill=orange!15, font=\small\bfseries},
    verify/.style={rectangle, draw, rounded corners, minimum width=2.8cm, minimum height=0.8cm, fill=purple!15, font=\small\bfseries},
    label/.style={font=\scriptsize, text=gray},
    arrow/.style={-{Stealth[length=2.5mm]}, thick},
    feedback/.style={-{Stealth[length=2mm]}, thick, gray!50}
]
    % === MAIN VERTICAL FLOW (center column) ===
    \node[phase] (preflight) {PREFLIGHT};
    \node[label, right=0.3cm of preflight] {Baseline assessment};

    \node[phase, below=1.2cm of preflight] (check) {CHECK};
    \node[label, right=0.3cm of check] {Sentinel gate};

    \node[decision, below=1.2cm of check] (proceed) {PROCEED?};

    % === BRANCHES ===
    \node[action, below left=1.2cm and 2cm of proceed, fill=green!10] (praxic) {PRAXIC};
    \node[label, below=0.05cm of praxic] {(execute)};

    \node[action, below right=1.2cm and 2cm of proceed, fill=blue!5] (noetic) {NOETIC};
    \node[label, below=0.05cm of noetic] {(investigate)};

    % === LOWER FLOW ===
    \node[phase, below=3.8cm of proceed] (postflight) {POSTFLIGHT};
    \node[label, left=0.3cm of postflight] {Learning delta};

    \node[verify, right=3cm of postflight] (posttest) {POST-TEST};
    \node[label, right=0.3cm of posttest] {Grounded verification};

    % === DUAL-TRACK UPDATES ===
    \node[update, below=1.2cm of postflight] (bayesian) {Bayesian Update};
    \node[label, below=0.05cm of bayesian] {Track 1: Self-referential};

    \node[update, below=1.2cm of posttest] (grounded) {Grounded Update};
    \node[label, below=0.05cm of grounded] {Track 2: Evidence-grounded};

    % === ARROWS (main flow) ===
    \draw[arrow] (preflight) -- (check);
    \draw[arrow] (check) -- (proceed);

    % Branch: YES -> PRAXIC, NO -> NOETIC
    \draw[arrow] (proceed) -- node[above left, font=\scriptsize] {YES} (praxic);
    \draw[arrow] (proceed) -- node[above right, font=\scriptsize] {NO} (noetic);

    % PRAXIC -> POSTFLIGHT (down and across, under the proceed diamond)
    \draw[arrow] (praxic.south) -- ++(0,-0.8) -| (postflight.north);

    % NOETIC -> CHECK (loop back via right side, outside all nodes)
    \draw[feedback] (noetic.east) -- ++(1.2,0) |- (check.east);

    % POSTFLIGHT -> POST-TEST
    \draw[arrow] (postflight) -- (posttest);

    % Dual-track updates
    \draw[arrow] (postflight) -- (bayesian);
    \draw[arrow] (posttest) -- (grounded);

    % Calibration feedback loop (to next transaction)
    \draw[feedback, dashed] (bayesian.west) -- ++(-2.5,0) |- node[pos=0.25, left, font=\scriptsize, text=gray] {next txn} (check.west);

\end{tikzpicture}
\caption{The 4-phase CASCADE workflow. PREFLIGHT establishes baseline, CHECK gates action, POSTFLIGHT measures learning delta, and POST-TEST grounds self-assessment in objective evidence. Dual-track calibration updates run in parallel. The NOETIC $\rightarrow$ CHECK feedback loop (gray solid) enables iterative investigation within a transaction. The calibration $\rightarrow$ CHECK feedback loop (gray dashed) shows how learned bias corrections inform future transactions' Sentinel decisions.}
\label{fig:cascade-workflow}
\end{figure}

\subsubsection{PREFLIGHT Phase}

At task onset, the AI assesses its epistemic state across all 13 vectors. This establishes the baseline against which learning will be measured.

\begin{lstlisting}[language=bash,caption={Example PREFLIGHT submission},showspaces=false,showstringspaces=false]
empirica preflight-submit - << EOF
{
  "session_id": "abc123",
  "vectors": {
    "know": 0.55,
    "uncertainty": 0.40,
    "context": 0.70,
    "engagement": 0.85
  },
  "reasoning": "Familiar domain but codebase unknown"
}
EOF
\end{lstlisting}

\subsubsection{CHECK Phase (Sentinel Gate)}

The Sentinel protocol evaluates whether the AI has sufficient knowledge to proceed. The gate applies bias corrections learned from historical calibration and returns:
\begin{itemize}
    \item \textbf{PROCEED}: Knowledge sufficient, enter praxic phase
    \item \textbf{INVESTIGATE}: Knowledge insufficient, remain in noetic phase
\end{itemize}

\subsubsection{POSTFLIGHT Phase}

After task completion (or significant progress), the AI reassesses its epistemic state. The delta between PREFLIGHT and POSTFLIGHT vectors measures learning:

\begin{equation}
\Delta_{\text{learning}} = V_{\text{POSTFLIGHT}} - V_{\text{PREFLIGHT}}
\end{equation}

This delta feeds the self-referential calibration track (Track~1), capturing systematic bias in initial self-assessment.

\subsubsection{POST-TEST Phase (Grounded Verification)}

After POSTFLIGHT, the framework automatically collects objective evidence to ground the self-assessment in measurable outcomes:

\begin{table}[h]
\centering
\begin{tabular}{llll}
\toprule
\textbf{Evidence Source} & \textbf{Quality} & \textbf{Vectors Grounded} \\
\midrule
Test results (pytest) & Objective & know, do, clarity \\
Git commit count \& diff stats & Objective & do, change, state \\
Git reflog (resets, abandoned branches) & Objective & uncertainty, signal \\
Git file churn (repeated edits) & Objective & uncertainty, coherence \\
Git diff granularity (word-diff) & Objective & know, clarity \\
Git commit timing patterns & Objective & do, state \\
Git diff-filter (add/modify/delete ratio) & Objective & state, completion \\
Goal/subtask completion & Semi-objective & completion, do, know \\
Artifact ratios (findings/dead-ends) & Semi-objective & know, uncertainty, signal \\
Issue resolution & Semi-objective & impact, signal \\
Sentinel decisions & Semi-objective & context, uncertainty \\
Commit message vs diff semantics & Semi-objective & clarity, signal \\
\bottomrule
\end{tabular}
\caption{Post-test evidence sources and the vectors they ground. Git-derived sources (6 of 12) provide Tier~1 objective evidence that the AI cannot influence.}
\label{tab:posttest-evidence}
\end{table}

The POST-TEST phase computes a calibration gap between self-assessed POSTFLIGHT vectors and evidence-derived estimates:

\begin{equation}
\text{gap}_{\text{vector}} = V_{\text{POSTFLIGHT}} - V_{\text{evidence}}
\end{equation}

Positive gaps indicate overestimation; negative gaps indicate underestimation relative to objective outcomes. This feeds the grounded calibration track (Track~2), which provides more trustworthy corrections because it anchors in what actually happened rather than another self-report.

Not all vectors can be grounded---ENGAGEMENT and DENSITY lack objective correlates. For these, self-referential calibration remains the only available signal.

\paragraph{Git as Epistemic Trajectory Store}

A significant fraction of post-test evidence derives from git, which passively records an epistemic trajectory parallel to the AI's self-report. Git's directed acyclic graph structure maps naturally to epistemic concepts: commits are discrete state transitions, branches are divergent investigations, merges are knowledge synthesis, and the reflog captures the true history including course corrections the AI may not self-report.

Six git-derived evidence sources are currently implemented. Of particular note:

\begin{itemize}
    \item \textbf{Reflog analysis}: The reflog captures every \texttt{reset}, \texttt{rebase}, and abandoned branch---objective evidence of dead-ends and backtracking that the AI has no mechanism to omit or distort. This is the strongest Tier~1 signal for UNCERTAINTY grounding.
    \item \textbf{File churn}: When the same file is edited $N$ times within a session, this indicates iteration and uncertainty about that component, regardless of what the AI's UNCERTAINTY self-assessment reports.
    \item \textbf{Commit timing}: Gaps between commits encode noetic-to-praxic transitions. Long gaps followed by rapid commits indicate investigation followed by execution---the workflow rhythm is detectable from timestamps alone.
    \item \textbf{Diff granularity}: Word-level diffs distinguish surgical edits (indicating understanding) from wholesale rewrites (indicating exploration), grounding KNOW without self-report.
\end{itemize}

Because git records accumulate passively during normal development, these evidence sources require no additional instrumentation---the epistemic trajectory is already being recorded.

\paragraph{Phase-Aware Evidence Collection}

The POST-TEST phase separates evidence collection into noetic and praxic phases when a CHECK boundary exists within the transaction:

\begin{itemize}
    \item \textbf{Noetic evidence} (PREFLIGHT $\rightarrow$ CHECK): Measures investigation quality---unknowns surfaced, dead-ends identified before action, findings logged during exploration, and Sentinel gate iterations. This grounds self-assessment of what the AI \emph{learned} during investigation.
    \item \textbf{Praxic evidence} (CHECK $\rightarrow$ POSTFLIGHT): Measures execution quality---goal completion, test results, git metrics, artifact ratios. This grounds self-assessment of what the AI \emph{accomplished} during action.
\end{itemize}

When no CHECK boundary exists (legacy sessions or single-phase transactions), evidence is collected in combined mode for backward compatibility. Phase-aware collection enables more precise calibration: an AI may be well-calibrated on investigation (noetic) but poorly calibrated on execution (praxic), or vice versa. Aggregate calibration would mask these patterns.

\paragraph{Scope-Weighted Artifact Metrics}

Not all artifacts contribute equally to calibration. Unknowns logged without connection to session goals (e.g., forward-looking research questions, general observations) are weighted at 0.3$\times$ when computing resolution ratios:

\begin{equation}
\text{resolution\_ratio} = \frac{\text{scoped\_resolved} + 0.3 \cdot \text{unscoped\_resolved}}{\text{scoped\_total} + 0.3 \cdot \text{unscoped\_total}}
\end{equation}

This prevents unscoped unknowns---which are deliberately left open for future investigation---from artificially depressing KNOW grounding. Logging unknowns \emph{is} epistemic work; leaving them intentionally unresolved should not penalize calibration. A floor of 0.3 is applied: even fully unresolved unknowns contribute base epistemic value (knowing what you don't know is itself knowledge).

\paragraph{Evidence Quality and Observation Variance}

Evidence sources have different reliability levels, encoded in observation variance during Bayesian updating:

\begin{table}[h]
\centering
\begin{tabular}{llr}
\toprule
\textbf{Quality} & \textbf{Sources} & \textbf{Obs.\ Variance} \\
\midrule
Objective & pytest, git metrics & 0.05 \\
Semi-objective & goal completion, artifacts, Sentinel & 0.07 \\
Inferred & derived ratios, indirect signals & 0.10 \\
Self-referential & POSTFLIGHT vectors (Track 1) & 0.10 \\
\bottomrule
\end{tabular}
\caption{Evidence quality hierarchy. Lower variance = higher trust = stronger pull on posterior. Objective evidence (tests, git) cannot be influenced by the AI and thus receives highest trust.}
\label{tab:evidence-quality}
\end{table}

The grounded calibration track (Track~2) uses observation variance $\sigma^2 = 0.05$ for objective evidence, compared to $\sigma^2 = 0.10$ for self-referential updates (Track~1). This gives objective evidence twice the ``pull'' on the posterior belief, anchoring calibration in what happened rather than what the AI reported.

\subsection{The Sentinel Protocol}

Sentinel is the epistemic gate that enforces noetic-before-praxic.

\subsubsection{Gate Logic}

\begin{algorithm}
\caption{Sentinel Decision}
\label{alg:sentinel}
\begin{algorithmic}[1]
\Procedure{SentinelDecision}{vectors, calibration}
    \State $\text{know}' \gets \text{vectors.know} + \text{calibration.know}$
    \State $\text{uncertainty}' \gets \text{vectors.uncertainty} + \text{calibration.uncertainty}$
    \If{$\text{know}' \geq 0.70$ \textbf{and} $\text{uncertainty}' \leq 0.35$}
        \State \Return PROCEED
    \Else
        \State \Return INVESTIGATE
    \EndIf
\EndProcedure
\end{algorithmic}
\end{algorithm}

\subsubsection{Calibration Integration}

The Sentinel applies corrections learned from historical Bayesian updates. If the AI systematically underestimates KNOW by 0.15, the gate adjusts:

\begin{itemize}
    \item Raw KNOW: 0.55
    \item Correction: +0.15
    \item Effective KNOW: 0.70
    \item Decision: PROCEED (would have been INVESTIGATE without correction)
\end{itemize}

This allows the system to compensate for learned biases while maintaining conservative defaults for uncalibrated vectors.

\subsection{Bayesian Belief Management}

The framework maintains Bayesian beliefs about each vector's calibration through two parallel tracks.

\subsubsection{Conjugate Normal Update}

Both tracks use the same conjugate normal inference:

\begin{equation}
\mu_{\text{posterior}} = \frac{\sigma^2_{\text{prior}} \cdot x + \sigma^2_{\text{obs}} \cdot \mu_{\text{prior}}}{\sigma^2_{\text{prior}} + \sigma^2_{\text{obs}}}
\end{equation}

\begin{equation}
\sigma^2_{\text{posterior}} = \frac{1}{\frac{1}{\sigma^2_{\text{prior}}} + \frac{1}{\sigma^2_{\text{obs}}}}
\end{equation}

The tracks differ in what constitutes an observation and in observation variance.

\subsubsection{Track 1: Self-Referential Calibration}

The self-referential track measures learning trajectory---how vectors change within an epistemic transaction:

\begin{equation}
\Delta_{\text{learning}} = V_{\text{POSTFLIGHT}} - V_{\text{PREFLIGHT}}
\end{equation}

\begin{itemize}
    \item \textbf{Observation}: POSTFLIGHT vector value
    \item \textbf{Observation variance}: $\sigma^2_{\text{obs}} = 0.10$ (self-reported, moderate uncertainty)
    \item \textbf{Measures}: Whether the AI consistently under- or overestimates at task onset
\end{itemize}

Positive delta indicates capability growth during the task. The mean delta becomes the bias correction applied at CHECK.

\subsubsection{Track 2: Grounded Calibration}

The grounded track measures calibration accuracy---whether POSTFLIGHT self-assessment matches objective outcomes:

\begin{equation}
\text{gap}_{\text{vector}} = V_{\text{POSTFLIGHT}} - V_{\text{evidence}}
\end{equation}

\begin{itemize}
    \item \textbf{Observation}: Evidence-derived vector estimate from POST-TEST
    \item \textbf{Observation variance}: $\sigma^2_{\text{obs}} = 0.05$ (evidence-based, lower uncertainty)
    \item \textbf{Measures}: Whether the AI's post-task self-assessment matches what actually happened
\end{itemize}

The lower observation variance gives grounded evidence stronger pull on the posterior, reflecting its higher trustworthiness compared to self-report. When the tracks diverge, grounded calibration takes precedence---it is anchored in outcomes, not self-assessment.

\subsubsection{Calibration Divergence}

The divergence between tracks reveals self-assessment quality:

\begin{equation}
\text{divergence}_{\text{vector}} = |\mu_{\text{self-ref}} - \mu_{\text{grounded}}|
\end{equation}

Low divergence indicates well-calibrated self-assessment. High divergence signals that the AI's self-reports are systematically disconnected from outcomes. Production data (\Cref{sec:results}) shows the largest divergence on IMPACT (+0.625 overestimation), with COMPLETION showing only +0.038 divergence after scope-aware evidence weighting. This suggests IMPACT self-assessment is particularly susceptible to overestimation.

\subsection{Data Model}

The framework persists state in SQLite with the following schema:

\paragraph{Sessions} Track AI interactions with unique identifiers, timestamps, and status.

\paragraph{Cascades} Complete PREFLIGHT $\rightarrow$ CHECK $\rightarrow$ POSTFLIGHT $\rightarrow$ POST-TEST loops with outcome tracking.

\paragraph{Epistemic Snapshots} Vector states at each CASCADE phase, enabling delta computation.

\paragraph{Bayesian Beliefs} Per-vector self-referential calibration state (Track~1): mean, variance, evidence count, and update history.

\paragraph{Grounded Beliefs} Per-vector grounded calibration state (Track~2): evidence-derived corrections with lower observation variance, updated from POST-TEST evidence.

\paragraph{Verification Evidence} Raw post-test evidence records per session, linking objective signals to vector estimates.

\paragraph{Assumptions} Unverified beliefs with confidence scores and urgency that increases with age. Urgency is computed as $\min(1.0,\; \frac{\text{age\_days}}{30} \times (1 - \text{confidence}))$, ensuring that low-confidence assumptions surface with increasing priority at each PREFLIGHT.

\paragraph{Decisions} Permanent audit-trail records of choice points: what was chosen, the rationale, reversibility classification (exploratory, committal, or forced), and optional outcome assessment with regret scoring.

\paragraph{Intent Edges} Provenance links from noetic artifacts to praxic outcomes, recording which findings and unknowns informed which goals and subtasks. These enable post-hoc analysis of what investigation led to what action.

\paragraph{Lessons} Executable procedures extracted from successful patterns: structured steps, prerequisites, corrections, and expected epistemic deltas (the vector changes a lesson should produce when applied). Lessons track replay history for reliability estimation and are subject to confidence decay via the cognitive immune system.

\subsection{CLI Interface}

Empirica provides a command-line interface for integration:

\begin{lstlisting}[language=bash,caption={Empirica CLI commands},showspaces=false,showstringspaces=false]
# Session management
empirica session-create --ai-id claude-code --output json
empirica session-list --status active

# CASCADE workflow (4-phase)
empirica preflight-submit -    # Baseline assessment
empirica check-submit -        # Sentinel gate: PROCEED or INVESTIGATE
empirica postflight-submit -   # Learning delta + auto POST-TEST

# Knowledge management
empirica finding-log --session-id ... --finding "..." --impact 0.8
empirica unknown-log --session-id ... --unknown "..."

# Dual-track calibration
empirica calibration-report --ai-id claude-code            # Track 1
empirica calibration-report --ai-id claude-code --grounded  # Track 2
empirica calibration-report --ai-id claude-code --trajectory # Trend
\end{lstlisting}

\subsection{Integration Patterns}

\subsubsection{IDE Integration}

Empirica integrates with development environments via hooks that trigger CASCADE phases at appropriate points in the development workflow.

\subsubsection{Agent Framework Integration}

For autonomous agents, Empirica provides epistemic checkpoints that gate action based on assessed knowledge state, implementing the noetic-before-praxic principle programmatically.

\subsection{Semantic Memory and Retrieval}

\subsubsection{Four-Layer Memory Architecture}

Empirica organizes memory across four layers with distinct latency and persistence characteristics:

\begin{table}[h]
\centering
\begin{tabular}{llll}
\toprule
\textbf{Layer} & \textbf{Storage} & \textbf{Latency} & \textbf{Persistence} \\
\midrule
HOT & In-context window & Nanosecond & Session-scoped \\
WARM & SQLite per-project & Microsecond & Permanent (structured) \\
SEARCH & Qdrant semantic index & Millisecond & Permanent (embeddings) \\
COLD & Git notes + YAML & Millisecond & Versioned, portable \\
\bottomrule
\end{tabular}
\caption{Four-layer memory architecture.}
\label{tab:memory-layers}
\end{table}

HOT memory is the AI's active context window---ephemeral and lost on compaction. WARM memory stores structured data (sessions, cascades, beliefs, goals, findings) in per-project SQLite databases. SEARCH memory provides semantic retrieval across four Qdrant collections: \emph{docs} (reference documentation), \emph{memory} (project artifacts), \emph{eidetic} (facts with confidence scores), and \emph{episodic} (session narratives with temporal decay). COLD memory archives artifacts in git notes and YAML files, making them version-controlled, portable, and accessible without a running database.

\subsubsection{Contextual Injection at PREFLIGHT}

At PREFLIGHT, before the AI has investigated the current task, the framework retrieves calibration-relevant context from the SEARCH layer. Retrieval is ranked by:

\begin{equation}
\text{rank} = \text{impact} \times \text{type\_confidence} \times \text{recency\_decay}
\end{equation}

\noindent where recency follows a 24-hour half-life, ensuring critical recent knowledge surfaces first while older context decays naturally. This serves four functions:

\begin{enumerate}
    \item \textbf{Pattern warnings}: Prior sessions with similar tasks and poor calibration scores surface as warnings (e.g., ``In session X with a similar task, you overestimated COMPLETION by +0.30'').
    \item \textbf{Dead-end avoidance}: Previously logged dead-ends matching the current task context appear at PREFLIGHT, preventing re-exploration of known failures.
    \item \textbf{Lesson retrieval}: Relevant procedural knowledge from prior sessions informs the baseline assessment.
    \item \textbf{Assumption surfacing}: Unresolved assumptions with high urgency appear at PREFLIGHT, ensuring unverified beliefs cannot be indefinitely deferred.
\end{enumerate}

POSTFLIGHT vectors, findings, and calibration outcomes are automatically embedded into the SEARCH layer, creating a growing corpus of episodic calibration memory. Cross-project global learnings can also be retrieved, enabling knowledge transfer between distinct projects. This enables the system to learn not just aggregate bias corrections (Bayesian beliefs) but also \emph{context-sensitive} corrections---recognizing that calibration accuracy varies by task type, domain, and complexity.

\subsubsection{Compaction Recovery}

When the AI's context window is compressed (a routine event in long-running sessions), the framework recovers context using epistemically weighted selection rather than chronological ordering. The recovery process detects whether the previous transaction completed (full loop) or was interrupted (mid-transaction), then reconstructs the most valuable context: high-impact findings, unresolved unknowns, active goals, and the transaction file needed for Sentinel continuity. This enables sessions that span compaction boundaries without catastrophic knowledge loss.

\subsubsection{Cognitive Immune System}

The lesson system includes a built-in challenge mechanism: when a new finding is logged, semantically similar lessons have their confidence reduced by 0.05 per match, with a floor at 0.3. This is domain-scoped---findings in one domain do not decay lessons in another. The metaphor is immunological: lessons encode successful procedures (antibodies); findings may invalidate them (antigens). The floor ensures lessons degrade gracefully rather than disappearing entirely, preserving potentially valuable procedures even when challenged.

\subsection{Production Deployment}

The framework has been deployed in production for 5 months (November 2025 -- March 2026):

\begin{table}[h]
\centering
\begin{tabular}{lr}
\toprule
\textbf{Metric} & \textbf{Value} \\
\midrule
Sessions completed & 792 \\
Evidence observations & 1,190,265 \\
Full CASCADE cycles & 616 \\
Clean learning pairs & 456 \\
Findings logged & 2,102 \\
Primary AI & Claude (Anthropic) \\
\bottomrule
\end{tabular}
\caption{Production deployment statistics.}
\label{tab:deployment-stats}
\end{table}

Cross-validation has been performed on GPT, Gemini, and Qwen models. The data presented in \Cref{sec:results} derives from this production deployment.

\subsection{Extended Subsystems}

Several subsystems extend the core framework to address practical deployment challenges.

\subsubsection{Skill Extraction}

Reference documentation for tools and frameworks consumes substantial context window space, much of which is irrelevant to the current task. Empirica addresses this through \emph{skill extraction}: verbose reference documents are distilled into structured skill cards that preserve decision-relevant knowledge (when to use, anti-patterns, cost models, key commands) while achieving 80--90\% token reduction. Skills are stored as markdown cards parsed into YAML runtime objects, and loaded into context only when semantically relevant to the current task.

\subsubsection{Emerged Personas}

Rather than prescribing fixed behavioral profiles, Empirica extracts \emph{emerged personas} from data: patterns of initial vectors, delta trajectories, and convergence thresholds observed in successful session branches. Each persona captures a working style---initial epistemic configuration, expected learning trajectory, and convergence threshold---with reputation tracking (uses count, success count). Personas are descriptive, not prescriptive: they represent observed modes of effective operation, not constraints on behavior.

\subsubsection{Attention Budget}

To prevent unbounded exploration, the framework implements an \emph{attention budget}: an information gain estimator using Shannon entropy of uncertainty, a knowledge gap multiplier, and diminishing returns decay. The budget is domain-scoped---investigation in one domain does not exhaust the budget for another. For multi-agent configurations, the budget gates subagent spawning: if the parent's attention budget is exhausted, a strong warning is issued before creating new subagents. Tool calls from subagents are added to the parent's budget consumption as delegated work.

\subsubsection{Blindspot Detection}

The grounded verification track reveals not just calibration errors but \emph{systematic blindspots}: dimensions the AI consistently fails to assess accurately. Blindspot detection performs negative space analysis on grounded calibration gaps---patterns of persistent divergence between self-assessed and evidence-derived vectors. Critical blindspots can override the CHECK gate: even if the AI's self-assessed vectors meet the PROCEED threshold, the Sentinel may force investigation if a historically ungrounded vector is critical to the current task.
